\appendix
\chapter{Release Verification}

The executable examples target serQ v0.1.0, commit
\texttt{f9fe9f2}, with IR version 9. The accompanying verifier runs
four FIFO service laws, three PS service laws, three finite-population
sizes, four PD bottleneck configurations, and three seeds of the
continuous tandem. It also recomputes the feedback example's roots and
checks the step-engine miss-price approximation.

\paragraph{Results that retain their assumptions.} PK, PS insensitivity,
finite-source MVA and the PD bottleneck formula agree with the simulations
within the verifier's stated finite-run tolerances. The feedback example
still has equilibria $0$, approximately $0.250$ and $0.882$ at
$\lambda=0.25$/s, and approximately $0.885$ at $\lambda=0.18$/s.
The step-engine test compares increased prefill occupancy with the FIFO
bracket at mean prefill availability. That is a check of an approximation,
not a proof that the engine is M/G/1. Its decode occupancy changes slightly
because prefill affects shared iteration time.

\paragraph{A changed result in the previous PD example.} The old PS
capacity expression used a name that also denoted the workload's appended
token count. The interpreter resolved it as that attribute, rather than
as the number at the station, giving almost every decode the full capped
throughput even when alone. Using \texttt{present} in
\texttt{ps(min(present,16))} restores the stated model. A turn with
mean decode work $200\times0.0002=0.04$ seconds cannot receive service in
about $0.0025$ seconds at a per-turn rate of at most one.
The table is generated from three matched seeds of the preceding example
and its serQ v0.1.0 port; it contains simulated means, not measurements.
\begin{center}\small
\begin{tabular}{@{}lrr@{}}
\toprule
Metric (seconds) & Previous runtime & serQ v0.1.1 \\
\midrule
TTFT (prefill instance) & 0.04327 & 0.04357 \\
Turn response & 0.05117 & 0.08868 \\
Decode service & 0.00249 & 0.03969 \\
\bottomrule
\end{tabular}
\end{center}

The prefill cost and TTFT remain close, while decode and whole-turn response
change. Matching a few seeds does not establish equality of all trajectories.

\paragraph{Results that need new interpretation.} A tail-block eviction has
a concave recomputation cost, so the fractional-knapsack proof does not
establish optimal block eviction. Partial reuse requires its own work
moments. The explicit KV-transfer operation reserves the destination and
releases the source after transfer; an ideal tandem that queues at the
destination only afterwards has different holding times. A preempted decode
retains generated tokens and re-prefills only lost KV. Waiting-prefix pinning
can protect one turn while taking capacity from others, so its net effect
requires the program and workload to be specified.

\paragraph{Reproduction.} At the repository root, run
\texttt{make lecture-results} for current-release checks.
\texttt{scripts/check\_lecture\_results.py} optionally accepts a preceding
runtime and its program directory to regenerate the comparison evidence
and this table. The release's pool, PD-transfer, preemption and recorded
vLLM-oracle tests check the underlying interpreter rules. The mathematical
proofs in these notes establish results for their stated models; those tests
do not turn every executable program into a formally proved queue.
